\section{Initial Kähler metrics on the full smoothing base}
\label{sec:initial-metrics}

We use the polarized family
$(X_p,\mathcal L_p,\Omega_p)$ of Section~\ref{sec:geometry}.
Let $\omega_0\in c_1(\mathcal L_0)$ be its central singular
Ricci-flat metric. The central conical theorem of
\cite[Theorem~1.4 and Lemma~A.1]{hein-sun} applies: the nodes
have strongly regular Calabi--Yau cones, and a full-support
line in the period base gives a polarized smoothing over a disc.
It supplies fixed holomorphic conical charts with potentials
\begin{equation}\label{eq:central-conical-potentials}
 \phi_{0,i}=\phi_{C,i}+v_i,\qquad
 \phi_{C,i}=c_i^{\,0}|z|^{4/3},\qquad
 |\nabla_C^jv_i|\le C_jr_i^{2+\mu-j}
\end{equation}
for some $\mu>0$. Here $r_i$ is the fixed cone radius,
comparable to $|z|^{2/3}$. We decrease $\mu$ so that $\mu\le3/2$.
The constant $c_i^{\,0}$ in the potential is unrelated to the
link norm $c_i$ in the capacity energy.

\begin{proposition}\label{prop:initial-kahler}
For every sufficiently small fixed $R>0$ and all sufficiently
small $p\in B^\circ$, there is a smooth Kähler form
$\widetilde\omega_p^R\in c_1(\mathcal L_p)$ with the following
properties.

On $r_i<R$ it is a normalized Stenzel form $\sigma_{i,p}$,
of length scale
\begin{equation}\label{eq:normalized-cap-scales}
 \varepsilon_i=b_i(p)^{1/2}|\tau_i(p)|^{1/3},
 \qquad b_i(p)\longrightarrow b_i(0)>0.
\end{equation}
Outside $r_i<2R$ it converges smoothly to $\omega_0$ under
exterior family identifications.

Let $\Theta_p$ be the positive density associated with
$\Omega_p\wedge\overline{\Omega_p}$, with the fixed conventional
complex factor, and put
\[
 V=\int_{X_p}c_1(\mathcal L_p)^3,\quad
 I(p)=\int_{X_p}\Theta_p,\quad
 \kappa(p)=V/I(p),\quad
 F_p^R=\log\frac{\kappa(p)\Theta_p}
                         {(\widetilde\omega_p^R)^3}.
\]
Then $I(p)\to I(0)>0$. There is a smooth radius function
$\varrho_p$, comparable to the scaled Stenzel radius on the
node regions and to one on the fixed exterior, such that, for
every fixed $k$,
\begin{equation}\label{eq:initial-residual}
 \limsup_{\substack{p\to0\\p\in B^\circ}}
 \max_{j\le k}\sup_{X_p}
       \varrho_p^j|\nabla^j F_p^R|\le C_kR^\mu.
\end{equation}
On each Stenzel region the stronger identity and estimate hold:
\begin{equation}\label{eq:pluriharmonic-residual}
 F_p^R=\log\left|\frac{a_i(z,p)}{a_i(0,p)}\right|^2,
 \qquad |\nabla^jF_p^R|\le C_j\varrho_p^{3/2-j},
 \qquad \Omega_p=a_i(z,p)\Omega_{Q,\tau_i}.
\end{equation}
The derivatives use the initial metric. The constants are
independent of the ratios and phases of the node parameters.
The parameter neighbourhood can depend on $R$ and $k$.
\end{proposition}

\subsection{Family charts and Hermitian weights}

The conical charts in \eqref{eq:central-conical-potentials} may
be chosen compatibly with nearby quadric fibres. Indeed lift
the fixed central transition to an invertible ambient holomorphic
germ $\Phi$. Since it preserves the node, $Q\circ\Phi=uQ$ for
a holomorphic unit $u$. After choosing a square root, the map
$\Psi=u^{-1/2}\Phi$ satisfies $Q\circ\Psi=Q$ and is invertible.
It therefore acts on all nearby fibres. The nonconstant scalar
factor changes the central metric at relative order
$O(|z|)=O(r_i^{3/2})$; its constant part changes only the fixed
normalization and phase. Thus
\eqref{eq:central-conical-potentials} holds in these family
charts after decreasing $\mu$ if needed.

The central form is the normalized Chern form of a Hermitian
metric $h_0$ on $\mathcal L_0$ over $X_0^{\mathrm{reg}}$:
use the bounded-potential Calabi--Yau solution relative to a
smooth background Chern form. We normalize Hermitian weights
so that their Chern form is $\mathrm i\partial\bar\partial\phi$
and represents $c_1$; the factor $2\pi$ is included in the
weight convention.

Choose a holomorphic frame of $\mathcal L$ on each node
neighbourhood. Its $h_0$-weight can be made exactly
$\phi_{0,i}$. The difference is pluriharmonic on the punctured
node. Since its link is $S^2\times S^3$, its first real
cohomology vanishes; the difference is the real part of a
single-valued holomorphic function. Hartogs extension on the
normal threefold node extends that function across the vertex.
Lift its germ to the ambient chart and extend the resulting
frame change to the family.

For fixed $R$, the complement of the radius-$R/2$ node
regions is a compact smooth part of the family. Extend $h_0$
smoothly there, for example by extending its logarithm relative
to a background metric in a smooth tubular neighbourhood and
using a partition of unity. Denote the extension by
$h_{\mathrm{ext},p}$. On each fixed collar its weight
$\phi_{\mathrm{ext},i,p}$ tends to $\phi_{0,i}$ in every
$C^k$ norm. No rate uniform as $R\downarrow0$ is required.

\subsection{Volume and the exact local normalization}

Choose a holomorphic nonvanishing ambient representative for
the coefficient $a_i$ in \eqref{eq:pluriharmonic-residual}.
Its value at $(0,p)$ is taken in this chosen representative,
even though $z=0$ does not belong to the smooth quadric.

We first prove continuity of $I(p)$ without restricting the
relative scales. Under $z=\sqrt{\tau}\,y$, the residue form
pulls back to $\tau\Omega_{Q,1}$ and the Stenzel metric scales
by $|\tau|^{2/3}$. The unit Stenzel volume has Euclidean
growth in real dimension six. Consequently the
$\Theta_{Q,\tau}$-volume of $r_i<s$ is at most $Cs^6$,
uniformly for small $\tau$. If the region meets the cap
core, the same estimate follows from the unit-cap volume
bound for radii above its positive minimum. Boundedness
of $a_i$ preserves this estimate. On the complement of
these balls the densities converge smoothly. Letting $s$
tend to zero proves $I(p)\to I(0)$ along every sequence in
the base. The limit is positive on the fixed exterior.

The intersection number $V$ is constant in the flat polarized
family, so $\kappa(p)\to\kappa(0)>0$. Normalize the holomorphic
scaling of the unit Stenzel form by
$\omega_{\mathrm{St},\tau}^3=\kappa_{\mathrm{St}}\Theta_{Q,\tau}$
and set
\begin{equation}\label{eq:stenzel-volume-normalization}
 b_i(p)=
 \left(\frac{\kappa(p)|a_i(0,p)|^2}
                  {\kappa_{\mathrm{St}}}\right)^{1/3},
 \qquad
 \sigma_{i,p}=b_i(p)\omega_{\mathrm{St},\tau_i(p)}.
\end{equation}
The limiting cone form of $\sigma_{i,p}$ is exactly
$\omega_{C,i}$: both have the standard quadric shape in
the chosen conical chart, and their volume coefficient is
$\kappa(0)|a_i(0,0)|^2$. That coefficient fixes the positive
metric scale. This proves \eqref{eq:normalized-cap-scales}.
The Stenzel potentials $\psi_{i,p}$ can be normalized by
constants so that they converge to $\phi_{C,i}$ with all
derivatives on each fixed punctured annulus.

\subsection{Gluing in the prescribed polarization}

Choose $\chi_i=0$ on $r_i\le R$ and $\chi_i=1$ on
$r_i\ge2R$, with $|\nabla^j\chi_i|\le C_jR^{-j}$.
In the node frame use the Hermitian weight
\begin{equation}\label{eq:initial-weight-gluing}
 \widetilde\phi_{i,p}
   =\psi_{i,p}+\chi_i
                  (\phi_{\mathrm{ext},i,p}-\psi_{i,p}).
\end{equation}
It agrees with the exterior weight near the outer overlap.
These weights define a global Hermitian metric on the
original line bundle $\mathcal L_p$. Its normalized Chern
form therefore represents exactly $c_1(\mathcal L_p)$.

The form is positive on the core because it is Stenzel,
and on the fixed exterior by convergence to $\omega_0$.
On the transition annulus the limit at fixed $R$ is
$\mathrm i\partial\bar\partial(\phi_{C,i}+\chi_iv_i)$.
Every differentiated cutoff term has relative size
$O(R^\mu)$ by \eqref{eq:central-conical-potentials}.
For $R$ small the limiting form is positive, and for
all sufficiently small $p$ positivity follows by smooth
convergence on this fixed annulus.

On the core, \eqref{eq:stenzel-volume-normalization} makes
the logarithmic residual exactly
\eqref{eq:pluriharmonic-residual}. A holomorphic logarithm
exists on a smaller ambient neighbourhood. Its Taylor
expansion has no $z$-constant term. A linear coordinate
has conic degree $3/2$; its differentiated estimates are
uniform on every scaled cap, including its core. The
higher Taylor terms have higher degree, and the convergent
series gives the stated estimate at every fixed order.
On the transition annulus the metric error is $O(R^\mu)$
with scaled derivatives, and the central volume coefficient
varies by $O(R^{3/2})$. On the fixed exterior the residual
tends smoothly to zero. These estimates prove
\eqref{eq:initial-residual} and complete the proof of
Proposition~\ref{prop:initial-kahler}. Notice the exact
compatibility
\begin{equation}\label{eq:initial-volume-compatibility}
 \int_{X_p}e^{F_p^R}(\widetilde\omega_p^R)^3=V.
\end{equation}

\subsection{Comparison with the surgery metrics}

The initial metrics satisfy
\begin{equation}\label{eq:initial-model-comparison}
 \limsup_{p\to0}\sup
 |\Phi_p^*\widetilde g_p^R-
       \bar g_{\varepsilon(p)}^{R,A}|
                  _{\bar g_{\varepsilon(p)}^{R,A}}
 \le C R^\mu+C e_{\mathrm{AC}}(A),
 \qquad e_{\mathrm{AC}}(A)\longrightarrow0.
\end{equation}
Here the model metrics and their fixed central and cap
modifications are those of Sections~\ref{sec:model-energy}
and \ref{sec:model-compactness}. On the fixed exterior
the initial metric tends to $g_0$, and its difference
from the central model is $O(R^\mu)$. On each inner region
it is exactly normalized Stenzel. The completed cap differs
from that metric only by the fixed large-radius
asymptotically conical error $e_{\mathrm{AC}}(A)$.
The same bound covers the intervening annuli.

The identifications can also accommodate all phases.
The normal graph map between $Q=\tau$ and $Q=0$ tends
to the identity away from the core; identifying the core
with $Q=1$ uses a choice of $\sqrt{\tau}$. The two cone-end
maps can differ by a bounded Reeb rotation, an isometry
of the cone. Interpolate that rotation in $\log r$
between $2A\varepsilon_i$ and $R$, keeping it constant
near the endpoints. Its radial angular derivative in
the unit frame is at most
$C/\log(R/(2A\varepsilon_i))$. Thus its relative metric
error tends to zero as that neck length tends to infinity.
The finite maximum tends to zero without a ratio bound.
The maps are chosen fibrewise; a single-valued
trivialization around the discriminant is not asserted.
This proves \eqref{eq:initial-model-comparison} on the
entire manifold, in the order: fixed $R,A$, then $p\to0$.

The construction parallels the Hermitian-weight method in
\cite[Proposition~5.9]{benabida}; the uniform estimates here
use the full period base. Neither
\eqref{eq:initial-residual} nor
\eqref{eq:initial-model-comparison} yet compares the
actual Ricci-flat metric with the surgery model.
